% =====================================================================
% Section: Conclusion -- SKELETON. The proposed reporting standard is the part
% a reader should be able to copy into a model card, so keep it short, concrete,
% and free of hedging.
% =====================================================================

\section{Conclusion: a reporting standard for compression claims}
\label{sec:conclusion}

The claim ``near-lossless'' is cheap to write and expensive to support. This
paper measured the gap---\textbf{4 of 12 determinate claims underpowered for
their own assertion, 5 of 17 unevaluable from what is reported, 0 of 17
reproducible from released artifacts}---and supplied the apparatus that closes
it: equivalence certification at a declared margin, and sample-size tables
computed from observed churn rather than assumed variance.

% Sequential certification was a full section in earlier drafts. It is a
% registered-but-unrun component and a results-free section reads as padding, so
% it is reduced to this sentence (2026-07-26). Restore the section only when the
% dated registration document exists under docs/ AND the component has run.
The same apparatus is being extended to anytime-valid sequential certification,
in which a confidence sequence lets an evaluation be monitored continuously and
stopped as soon as the model is certified; that component is registered but has
not yet been run, and no results for it are reported here.

We propose five lines that a model card or method paper can adopt directly.

\begin{enumerate}
  \item \textbf{Declare a margin.} ``Equivalent'' without $\pm m$ is not a
  testable statement. State the margin before evaluating.
  \item \textbf{Run the paired test.} Report TOST at your margin, not just a
  non-significant difference test. Failing to detect a difference is not
  equivalence.
  \item \textbf{Report churn next to net delta.} They are different quantities;
  in the public record churn runs five to six times the net delta, and roughly
  one evaluation cell in ten posts an identical score while still disagreeing on
  items.
  \item \textbf{Cite the sample size you met.} Table~\ref{tab:certification}
  gives the count your benchmark family requires at your margin; say which
  column you are in and what you actually ran.
  \item \textbf{Release per-item outputs.} A few megabytes converts an assertion
  into something a third party can check. No source we audited did this for the
  tasks its own claim covers.
\end{enumerate}

\TODO{closing paragraph on the H3 experiment's contribution to the standard,
written only after all registered cells exist; it must not state an H3 outcome
until then}
